% =====================================================================
\section{LRL by Example}\label{sec:tour}
% =====================================================================

We introduce LRL with the running example of this paper, case study~B (\S\ref{sec:eval:cases}): a channel that may carry exactly as many messages as a vector has elements.
The complete program is in Appendix~\ref{app:listings}; it compiles and runs with both of LRL's code generators.

\paragraph{Syntax.}
LRL programs are S-expressions.
\lstinline|(def name Type value)| defines a constant; \lstinline|(pi x A B)| is the dependent function type $\Pi(x{:}A).\,B$ and \lstinline|(lam x A body)| a function; braces mark implicit arguments, which the elaborator infers.
\lstinline|(sort 0)| is the universe of propositions and \lstinline|(sort 1)| the universe of ordinary types.
\lstinline|(inductive Name Type (ctor c T) ...)| declares an inductive type with constructors.

\paragraph{Vectors.}
A vector of length $n$ is an inductive family indexed by a natural number:
\begin{lstlisting}
(inductive Vec (pi A (sort 1) (pi n Nat (sort 1)))
  (ctor vnil  (pi {A (sort 1)} (Vec A zero)))
  (ctor vcons (pi {A (sort 1)} (pi {n Nat}
                (pi h A (pi t (Vec A n) (Vec A (succ n))))))))
\end{lstlisting}

\paragraph{An owned channel.}
A channel of type \lstinline|(Chan n)| may still send \lstinline|n| messages; it records what it has sent in a list.
The marker \lstinline|affine| says that a channel is a resource: it cannot be copied, so a channel value can be used at most once.
\begin{lstlisting}
(inductive (affine) Chan (pi n Nat (sort 1))
  (ctor mk_chan (pi {n Nat} (pi log (List Nat) (Chan n)))))

(def close (pi c (Chan zero) Receipt)
  (lam c (Chan zero) (match c Receipt (case (mk_chan log) (receipt log)))))
\end{lstlisting}
\lstinline|close| accepts only a channel on which no message remains to be sent.

\paragraph{Operations generated by a macro.}
Sending is the same for every payload type, so the operations are generated by a macro.
\lstinline|(defsend NAME PAYLOAD ENCODE)| defines an operation that takes a \lstinline|(Chan (succ n))|, consumes it, and returns a \lstinline|(Chan n)| (Listing~\ref{lst:defsend}).
\begin{lstlisting}[float=*,caption={The \texttt{defsend} macro and the two calls that generate \texttt{send} and \texttt{send\_bit} (\texttt{case\_studies/lrl/protocol.lrl}).},label={lst:defsend}]
(defmacro defsend (name payload encode)
  (def name (pi {n Nat} (pi c (Chan (succ n)) (pi #[once] x payload (Chan n))))
    (lam {n} Nat (lam c (Chan (succ n)) (lam #[once] x payload
      (match c (Chan n)
        (case (mk_chan log) (mk_chan (cons (encode x) log)))))))))

(def log (pi m Nat Nat) (lam m Nat (print_nat m)))   ; prints a message, returns it
(defsend send Nat log)
(defsend send_bit Bool bit)
\end{lstlisting}
The annotation \lstinline|#[once]| is a function kind (\S\ref{sec:core:kinds}): after \lstinline|send| has received its channel, the remaining function owns that channel, so it may be called only once.
The macro is hygienic: the template binds a variable named \lstinline|log|, and the call \lstinline|(defsend send Nat log)| passes the user's function \lstinline|log|; the two do not interfere (\S\ref{sec:macros:hygiene}).

\paragraph{All three together.}
\lstinline|send_all| (Listing~\ref{lst:sendall}) sends every element of a vector of length \lstinline|n| over a channel that expects exactly \lstinline|n| messages, and closes it.
\begin{lstlisting}[float=*,caption={Sending a vector over a channel: dependent types, ownership and macro-generated operations in one function.},label={lst:sendall}]
(def send_all (pi {n Nat} (pi v (Vec Nat n) (pi #[once] c (Chan n) Receipt)))
  (lam {n} Nat (lam v (Vec Nat n)
    (match v (motive (lam k Nat (lam w (Vec Nat k) (pi #[once] c (Chan k) Receipt))))
      (case (vnil) close)
      (case (vcons h t ih) (lam #[once] c (Chan (succ _)) (ih (send c h))))))))
\end{lstlisting}
The \lstinline|match| has an explicit \emph{motive}: its result type depends on the length \lstinline|k| of the vector being matched, so the empty case must produce a function from \lstinline|(Chan zero)| (that is \lstinline|close|), and the case for \lstinline|(vcons h t)| must produce a function from \lstinline|(Chan (succ k))|.
That function sends the head and passes the channel to \lstinline|ih|, the result of the recursive call on the tail.
The type ties the data to the protocol, the kind ties each channel to a single use, and the operations come from a macro.
Because LRL has no way to hide constructors, the guarantee holds for code that uses channels through these operations; code that applies \lstinline|mk_chan| directly can create a channel at any index (\S\ref{sec:eval:cases}).

\paragraph{What the compiler rejects.}
Each of the following mistakes is a variant of the program in the repository, and is rejected before any code is generated:
\begin{itemize}\setlength\itemsep{1pt}
\item sending twice on the same channel value: the kernel reports \texttt{K0021} \lstinline|variable 'c' is used after it was moved|;
\item the same mistake inside a macro: the same error, labelled with the macro call;
\item calling \lstinline|close| while a message remains, or calling \lstinline|send_all| with a vector of the wrong length: the elaborator reports a type error (\texttt{F0214}), because \lstinline|(succ zero)| is not \lstinline|zero|;
\item using the channel after a branch of a \lstinline|match| has consumed it: \texttt{K0021}.
\end{itemize}
Consuming the channel in \emph{both} branches of a match on a Boolean is accepted, because exactly one branch runs.
The discipline is affine, not linear: dropping a channel without closing it is accepted, as in Rust.

\paragraph{Proofs.}
Proofs are ordinary definitions whose types are propositions.
Case study~A (\S\ref{sec:eval:cases}) proves, for instance, that reversing a vector twice gives the original vector (Listing~\ref{lst:involution}), by induction on the vector, using a lemma about \lstinline|vsnoc| and the equality lemmas \lstinline|trans| and \lstinline|cong|, all defined in the same file (Appendix~\ref{app:listings}).
\begin{lstlisting}[float=*,caption={The involution of \texttt{vreverse}, a kernel-checked proof (\texttt{case\_studies/lrl/vectors.lrl}).},label={lst:involution}]
(def vreverse_involutive (pi {A (sort 1)} (pi {n Nat} (pi v (Vec A n)
                           (Eq (Vec A n) (vreverse (vreverse v)) v))))
  (lam {A} (sort 1) (lam {n} Nat (lam v (Vec A n)
    (match v (motive (lam k Nat (lam w (Vec A k) (Eq (Vec A k) (vreverse (vreverse w)) w))))
      (case (vnil) (refl (Vec A zero) vnil))
      (case (vcons h t ih) (trans (vreverse_vsnoc (vreverse t) h) (cong (vcons h) ih))))))))
\end{lstlisting}
The kernel checks such proofs like every other definition, and code generation erases them.
